\paragraph{Two certified numerical lower bounds.}
Using the first \(1000\) positive simple critical zeros and their
conjugates in \eqref{eq:fejer-critical-poles} gives a right side
strictly above \(45769/100000\). The calculation uses
outward-rounded ball arithmetic at \(128\) bits. The supplied
verification script isolates the zeros, checks that the ordinate
enclosures are disjoint, and encloses the derivative and residue
at each zero. The \texttt{python-flint}~0.9.0 interface calls
the rigorous FLINT zeta routines described in its
\href{https://flintlib.org/doc/acb_dirichlet/index.html\#zeros-of-the-riemann-zeta-function}%
{documentation}. No zero outside the stated set enters the sum.

A positive weighting in the logarithmic variable yields
the stronger certified bound
\begin{equation}\label{eq:fejer-weight-bound}
 \limsup_{N\to\infty}\frac{Q_\lambda(N)}{N}
 >\frac{3257}{5000}=0.6514.
\end{equation}
To prove it, let \(C\) denote the limit superior in
\eqref{eq:fejer-weight-bound}. If \(C<\infty\), the proof of
Proposition~\ref{prop:fejer-multipole} gives, for
\(F(t)=e^{-t/2}U_{e^t}\),
\(\limsup_{t\to\infty}\|F(t)\|_{\mathcal H}=\sqrt C\).
Let \(\rho_1=1/2+i\gamma_1\) be the first positive critical zero,
\(A_1=\zeta(2\rho_1)/(\rho_1\zeta'(\rho_1))\), and
\(a=-1/\zeta(1/2)>0\). For \(m=100\), put
\[
 P(t)=\frac{1}{m+1}
     \left|\sum_{j=0}^{m}e^{ij(\gamma_1t-\theta)}\right|^2,
 \qquad e^{i\theta}A_1=-|A_1|.
\]
The trigonometric polynomial is nonnegative and its constant
coefficient is one, so
\(\delta\int_0^\infty e^{-\delta t}P(t)\dd t\to1\).
For each \(\varepsilon>0\), the norm of \(F(t)\) is at most
\(\sqrt C+\varepsilon\) outside a fixed bounded interval. The
integral on that interval vanishes after multiplication by
\(\delta\). Thus the limit of
\(\delta\int_0^\infty e^{-\delta t}P(t)F(t)\dd t\), when it
exists, has norm at most \(\sqrt C\).

For \(2\le k\le100\), ball arithmetic verifies
\(\zeta(1/2+ik\gamma_1)\ne0\) by enclosing its value over the
whole ordinate enclosure. Equation~\eqref{eq:fejer-vector-mellin}
then shows that the terms at those frequencies and their
conjugates vanish in the limit. The constant term contributes
\(-a w_{1/2}\). The two terms at frequencies \(\pm\gamma_1\)
give \(-c\R w_{\rho_1}\) with
\(c=2m|A_1|/(m+1)\), since the first Fourier coefficient of
\(P\) has modulus \(m/(m+1)\). The limit therefore exists and
\[
 C\ge \|a w_{1/2}+c\R w_{\rho_1}\|_{\mathcal H}^2.
\]
Writing \(v=\R w_{\rho_1}\), direct integration gives
\begin{align*}
 \|w_{1/2}\|_{\mathcal H}^2&=\frac23,\\
 \|v\|_{\mathcal H}^2&=\frac32+\R\frac1{1+2\rho_1}
                       -2\R\frac1{1+\rho_1},\\
 \langle w_{1/2},v\rangle_{\mathcal H}
    &=\frac13-\R\frac1{1+\rho_1}+2\R\frac1{3/2+\rho_1}.
\end{align*}
The same script encloses the squared norm and verifies that
its lower endpoint exceeds the exact rational \(3257/5000\).
This proves \eqref{eq:fejer-weight-bound}, including the
immediate case \(C=\infty\).
